% DRAFT, rewrite
\section{Formal analysis using Alloy}
The rules about what a section may show are easy to break when parts arrive in any
order, so they are modelled in Alloy 6 (\texttt{docs/alloy/loading.als}). The model
has the states of a section from section 2.1.2, the events that move it (time
passing, an answer or a failure, retry, a later refresh, sign-out) and a value
computed from several parts.

\subsection{What is checked}
\begin{center}\small
\begin{tabularx}{\linewidth}{@{}l l X l@{}}
  \toprule
  Command & Scope & What it shows & Result \\
  \midrule
  \texttt{run firstSignIn} & 4, 10 steps & Three sections with no stored copy can all arrive & to be run \\
  \texttt{run slowThenRetried} & 4, 12 steps & A section can wait long, fail, be retried and arrive & to be run \\
  \texttt{run averageAppears} & 4, 10 steps & A value computed from two parts can appear & to be run \\
  \texttt{check EmptyOnlyAfterAnswer} & 4, 12 steps & \idref{R2}: never Empty before the part answered & to be run \\
  \texttt{check ComputedNeedsAllInputs} & 4, 12 steps & \idref{R3}: a computed value is visible only with all its parts & to be run \\
  \texttt{check ContentNeverLost} & 4, 12 steps & \idref{R6}, \idref{R17}: a failure never removes content & to be run \\
  \texttt{check CopyBelongsToSession} & 4, 12 steps & \idref{R14}, \idref{R40}: no copy of another account survives sign-out & to be run \\
  \texttt{check HeldShowsOwnCopy} & 4, 12 steps & A held section shows a copy of the account signed in & to be run \\
  \bottomrule
\end{tabularx}
\end{center}
% TODO: execute every command in the Alloy Analyzer and replace "to be run" with
% the result (instance found / no counterexample / counterexample).
The commands have not been executed yet; the results column is filled after a run
in the Alloy Analyzer.

\subsection{Model}
\lstinputlisting[language=Alloy]{../alloy/loading.als}
